% Inputs 1,2 are observed; 3,4 have indistinguishable opposing label laws.
% Volume three architecture figures. Independent export only; not book fonts.
\colorlet{MutedBlueFill}{FigureBlueFill}
\colorlet{MutedGreenFill}{FigureGreenFill}
\colorlet{MutedPurple}{FigurePurple}
\colorlet{MutedPurpleFill}{FigurePurpleFill}
\tikzset{
 v3 picture/.style={x=1mm,y=1mm,font=\figurefont\fontsize{8}{10}\selectfont,text=FigureInk,line width=.65pt},
 v3 box/.style={draw=FigureBlue,fill=white,rounded corners=1.8pt,line width=.7pt,inner sep=2.5pt,align=center},
 v3 data/.style={v3 box,fill=FigureBlueFill},
 v3 model/.style={v3 box,draw=FigureGreen,fill=FigureGreenFill},
 v3 output/.style={v3 box,draw=FigureYellow,fill=FigureYellowFill},
 v3 risk/.style={v3 box,draw=FigureRed,fill=FigureRedFill},
 v3 group/.style={draw=FigureBlue,rounded corners=1.8pt,line width=.65pt,inner sep=3mm},
 v3 flow/.style={draw=FigureStroke,line width=.75pt,-{Stealth[length=1.65mm,width=1.05mm]},rounded corners=1pt},
 v3 guide/.style={draw=FigureMuted,line width=.6pt,dashed},
 v3 feedback/.style={v3 flow,draw=FigureGreen,dashed},
 v3 title/.style={draw=none,fill=none,font=\figurefont\fontsize{8}{10}\selectfont,align=center,inner sep=1pt},
 v3 note/.style={draw=none,fill=none,font=\FigureNoteFont\fontsize{7.5}{9.5}\selectfont,text=FigureMuted,align=center,inner sep=1pt},
 v3 label/.style={draw=none,fill=none,font=\figurefont\fontsize{8}{10}\selectfont,align=center,inner sep=1pt},
 v3 wire/.style={draw=FigureMuted,line width=.45pt},
 v3 point/.style={circle,draw=FigureBlue,fill=FigureBlueFill,minimum size=2.3mm,inner sep=0pt,line width=.55pt}
}
% Bring pre-existing native TikZ geometry to the same typography and scale.
\tikzset{
 foundation picture/.append style={font=\figurefont\fontsize{8}{10}\selectfont,line width=.65pt},
 foundation node/.append style={rounded corners=1.8pt,line width=.7pt,font=\figurefont\fontsize{8}{10}\selectfont},
 foundation heading/.append style={draw=none,fill=none,font=\figurefont\fontsize{8}{10}\selectfont},
 foundation note/.append style={draw=none,fill=none,font=\FigureNoteFont\fontsize{7.5}{9.5}\selectfont},
 foundation flow/.append style={line width=.75pt,-{Stealth[length=1.65mm,width=1.05mm]}},
 foundation frame/.append style={rounded corners=1.8pt,line width=.7pt},
 foundation edge label/.append style={draw=none,fill=none,font=\figurefont\fontsize{8}{10}\selectfont}
}
% Official Lucide outline icons, consistently used across diagrams. Size in mm.
\newcommand{\VIcon}[5]{\node[inner sep=0pt] at (#2,#3) {\includegraphics[width=#4mm]{figures/03-foundations/00-shared/tikz/icons/#1-#5.pdf}};}
% A neutral network glyph (not a claim about an actual model architecture).
% Coordinates denote its centre; scale 1 gives 13 x 8 mm.
\newcommand{\VNet}[4]{\begin{scope}[shift={(#1,#2)},scale=#3]
 \foreach \a in {-3,0,3}{\foreach \b in {-3,0,3}{\draw[v3 wire] (-5,\a)--(0,\b);}}
 \foreach \a in {-3,0,3}{\draw[v3 wire] (0,\a)--(5,0);}
 \foreach \a in {-3,0,3}{\draw[draw=#4,fill=#4,line width=.55pt] (-5,\a) circle(.9);\draw[draw=#4,fill=white,line width=.55pt] (0,\a) circle(.9);}
 \draw[draw=#4,fill=#4,line width=.55pt] (5,0) circle(.9);
\end{scope}}
\newcommand{\VDots}[4]{\foreach \k in {0,1,2,3}{\fill[#4] ({#1+2.7*\k},#2) circle(#3);}}
% Feature vector/tensor elements are scientific glyphs, not decorative icons.
\newcommand{\VVector}[4]{\foreach \k in {0,...,4}{\fill[#4,rounded corners=.25pt] ({#1+2.2*\k},#2) rectangle ++(1.5,#3);}}

\begin{tikzpicture}[v3 picture]
\path[use as bounding box] (0,0) rectangle (163,71);
\node[v3 title] at (20,52) {共同训练样本 $S$};
\draw[v3 group] (2,24) rectangle (38,48);
\VIcon{bird}{11}{40}{7}{blue}\VIcon{car}{29}{40}{7}{blue}
\node[v3 label] at (11,30) {$(1,0)$};\node[v3 label] at (29,30) {$(2,1)$};
\node[v3 title] at (61,52) {同一算法 $A$};
\node[v3 model,minimum width=26mm,minimum height=16mm] (A) at (61,38) {};
\VNet{61}{38}{1}{FigureGreen}
\node[v3 label,anchor=north] at (61,27) {固定预测 $h(3)=0$\\$h(4)=1$};
\draw[v3 flow] (38,38)--(A.west);
% The two sources keep the observed pairs identical.
\foreach \y/\p/\col in {51/0/FigureBlue,19/1/FigurePurple}{
 \draw[v3 group,draw=\col] (87,{\y-13}) rectangle (162,{\y+14});
 \node[v3 title,text=\col] at (101,{\y+11}) {来源 $\mathcal P_{\p}$};
 \draw[v3 guide,draw=\col] (89,{\y-10}) rectangle (114,{\y+7});
 \VIcon{bird}{95}{\y+2}{5.5}{blue}\VIcon{car}{108}{\y+2}{5.5}{blue}
 \node[v3 label] at (95,{\y-6}) {$(1,0)$};\node[v3 label] at (108,{\y-6}) {$(2,1)$};
 \draw[v3 guide] (117,{\y-10})--(117,{\y+11});
 \VIcon{dog}{137}{\y+8}{5.5}{blue}\VIcon{ship}{153}{\y+8}{5.5}{blue}
 \node[v3 label] at (124,{\y+1}) {真值};
 \node[v3 label] at (124,{\y-7}) {预测};
 \node[v3 label] at (137,{\y+1}) {$\p$};\node[v3 label] at (153,{\y+1}) {$\p$};
}
\node[v3 label,text=FigureGreen] at (137,44) {$0\;\checkmark$};
\node[v3 label,text=FigureRed] at (153,44) {$1\;\times$};
\node[v3 label,text=FigureRed] at (137,12) {$0\;\times$};
\node[v3 label,text=FigureGreen] at (153,12) {$1\;\checkmark$};
\draw[v3 flow] (A.east)--(80,38)--(80,51)--(87,51);
\draw[v3 flow] (80,38)--(80,19)--(87,19);
\node[v3 note] at (125,2) {四个输入等概率，各为 $1/4$};
\end{tikzpicture}
